\documentclass[10.5pt,a4paper]{article}
% Estilo visual Ryoku: papel y tinta (PAPER #000000, INK #cdc4ba), sin sombras, radio minimo
\usepackage[utf8]{inputenc}
\usepackage[T1]{fontenc}
\usepackage{XCharter}          % titulares (alternativa serif a Fraunces en pdflatex)
\usepackage[sfdefault,scaled=0.95]{FiraSans}  % UI y cuerpo (Inter y Space Grotesk no estan en pdflatex)
\usepackage[scaled=0.9]{inconsolata}   % datos y etiquetas (alternativa a Space Mono)
\renewcommand{\familydefault}{\sfdefault}
\usepackage[spanish]{babel}
\usepackage[margin=2.3cm,top=2.2cm,bottom=2.4cm]{geometry}
\usepackage[table]{xcolor}
\usepackage{microtype,booktabs,tabularx,longtable,tikz,titlesec,fancyhdr,enumitem}
\usepackage[most]{tcolorbox}
\usepackage{needspace}
\newcommand{\cab}[1]{\mbox{\mono\footnotesize\color{inkmuted}#1}}
\definecolor{paper}{HTML}{000000}
\definecolor{ink}{HTML}{CDC4BA}
\definecolor{inkdim}{HTML}{B0A9A0}
\definecolor{inkmuted}{HTML}{958F87}
\definecolor{tint}{HTML}{121110}
\definecolor{tinthi}{HTML}{1E1C1A}
\definecolor{inkline}{HTML}{3A3733}   % ink al ~.26 sobre negro
\pagecolor{paper}\color{ink}
\usepackage[colorlinks=true,linkcolor=ink,urlcolor=inkdim,citecolor=ink]{hyperref}

\newcommand{\serif}{\rmfamily}
\newcommand{\mono}{\ttfamily}
\newcommand{\eyebrow}[1]{{\mono\footnotesize\color{inkmuted}\MakeUppercase{#1}}}
\newcommand{\cSi}{\textcolor{ink}{\large$\bullet$}}\newcommand{\cDudo}{\textcolor{inkmuted}{?}}\newcommand{\cNada}{\textcolor{inkmuted}{--}}
\newcommand{\nv}{\textcolor{inkmuted}{\textit{No verificado}}}

\titleformat{\section}{\serif\huge\color{ink}}{\mono\small\color{inkmuted}\twodigits{\thesection}\ //\ }{0pt}{}[\vspace{2pt}{\color{inkline}\titlerule}]
\titleformat{\subsection}[hang]{\serif\fontsize{28}{30}\selectfont\color{ink}}{\mono\small\color{inkmuted}\thesubsection\ //\ }{0pt}{}
\titleformat{name=\subsection,numberless}{\mono\small\color{inkdim}}{}{0pt}{}\titlespacing*{\subsection}{0pt}{16pt}{6pt}
\titlespacing*{\section}{0pt}{22pt}{10pt}
\titlespacing*{\subsection}{0pt}{16pt}{6pt}
\newcommand{\twodigits}[1]{\ifnum#1<10 0#1\else#1\fi}

\pagestyle{fancy}\fancyhf{}
\renewcommand{\headrulewidth}{0pt}
\fancyfoot[L]{\mono\footnotesize\color{inkmuted}RYOKU \textperiodcentered\ \docid}
\fancyfoot[R]{\mono\footnotesize\color{inkmuted}\thepage}
\fancypagestyle{plain}{\fancyhf{}\fancyfoot[R]{\mono\footnotesize\color{inkmuted}\thepage}}
\arrayrulecolor{inkline}
\setlength{\aboverulesep}{3pt}\setlength{\belowrulesep}{3pt}
\setlength{\parindent}{0pt}\setlength{\parskip}{5pt}
\setlist[itemize]{leftmargin=1.2em,itemsep=3pt,label={\color{inkmuted}--}}

\newcommand{\etiq}[1]{\tikz[baseline=(t.base)]\node[draw=inkmuted,line width=0.4pt,rounded corners=2pt,inner xsep=6pt,inner ysep=3pt,text=ink,font=\mono\footnotesize] (t) {\MakeUppercase{#1}};}
\newcommand{\kpi}[2]{\begin{minipage}[t]{0.235\textwidth}{\mono\scriptsize\color{inkmuted}\MakeUppercase{#1}}\\[2pt]{\serif\large\color{ink}#2}\end{minipage}}
\newenvironment{ficha}[6]{\needspace{22\baselineskip}\subsection{#1}\par\nopagebreak\vspace{-4pt}\etiq{#2}\par\vspace{8pt}%
 \begin{tcolorbox}[enhanced,colback=tint,colframe=inkline,boxrule=0.4pt,arc=2pt,left=8pt,right=8pt,top=7pt,bottom=7pt,boxsep=0pt]
 \kpi{Empleados}{#3}\hfill\kpi{Fundaci\'on}{#4}\hfill\kpi{Modelo}{#5}\hfill\kpi{Cloud}{#6}
 \end{tcolorbox}\vspace{6pt}{\color{inkline}\hrule height 0.4pt}\vspace{3pt}}{}
\newcommand{\campo}[2]{\noindent\hangindent=3.4cm\hangafter=1\makebox[3.4cm][l]{\mono\footnotesize\color{inkmuted}\MakeUppercase{#1}}#2\par\vspace{2pt}{\color{inkline}\hrule height 0.4pt}\vspace{3pt}}

\newcommand{\cPar}{\textcolor{inkmuted}{$\circ$}}
\newcommand{\crit}{\textcolor{ink}{\large$\bullet$}}
\newcommand{\perfilkpi}[2]{\begin{minipage}[t]{0.235\textwidth}{\mono\scriptsize\color{inkmuted}\MakeUppercase{#1}}\\[2pt]{\serif\large\color{ink}#2}\end{minipage}}

\newcommand{\docid}{Caso pr\'actico 1}
\usepackage{graphicx,fancyvrb}
\newcommand{\fig}[3][\textwidth]{\begin{center}\begin{tcolorbox}[colback=tint,colframe=inkline,boxrule=0.4pt,arc=2pt,left=3pt,right=3pt,top=3pt,bottom=3pt,width=#1,halign=center,nobeforeafter]\includegraphics[width=\linewidth]{capturas/#2.png}\end{tcolorbox}\\[3pt]{\mono\footnotesize\color{inkmuted}#3}\end{center}}
\usepackage{listings}
\definecolor{cgold}{HTML}{E6B84A}
\let\oldtexttt\texttt
\renewcommand{\texttt}[1]{\textcolor{cgold}{\oldtexttt{#1}}}
\definecolor{termbg}{HTML}{0B0D10}
\definecolor{termline}{HTML}{6B665F}
\definecolor{cprompt}{HTML}{7FD962}
\definecolor{ccmd}{HTML}{82AAFF}
\definecolor{cflag}{HTML}{FFB86C}
\definecolor{cstr}{HTML}{A5E075}
\definecolor{cvar}{HTML}{C792EA}
\definecolor{ccom}{HTML}{6B7280}
\definecolor{ctext}{HTML}{E6E1DA}
% comando en linea: caja con borde y color de terminal
\newcommand{\cmd}[1]{\tcbox[on line,colback=termbg,colframe=termline,boxrule=0.4pt,arc=2pt,left=2pt,right=2pt,top=1pt,bottom=1pt,boxsep=0pt,nobeforeafter]{{\mono\small\color{cprompt}\$\,\color{ctext}#1}}}
\lstdefinelanguage{shellx}{alsoletter={-},
  morekeywords=[1]{sudo,apt,systemctl,usermod,newgrp,docker,curl,nano,mkdir,cd,ip,psql,echo},
  morekeywords=[2]{-y,-aG,-d,-U,-c,-I,-br,-a,--rm,--now,enable,update,upgrade,install,run,compose,exec,up,ps,version,manifest,inspect,docker-compose-v2,docker.io},
  morecomment=[l]{\#},morestring=[b]",morestring=[b]'}
\lstdefinelanguage{yamlx}{alsoletter={-_},morekeywords=[1]{services,volumes,environment,ports,depends_on},
  morekeywords=[2]{db,web,image,container_name,restart,command,POSTGRES_DB,POSTGRES_USER,POSTGRES_PASSWORD,PGDATA,HOST,USER,PASSWORD},
  morecomment=[l]{\#},morestring=[b]"}
\lstdefinestyle{term}{basicstyle=\mono\small,keywordstyle=[1]\color{ccmd},keywordstyle=[2]\color{cflag},
  stringstyle=\color{cstr},commentstyle=\color{ccom}\itshape,showstringspaces=false,columns=fullflexible,keepspaces=true,
  literate={\$\ }{{\color{cprompt}\$\ }}2 {\$USER}{{\color{cvar}\$USER}}5,upquote=true}
\tcbset{termbox/.style={enhanced,breakable,colupper=ctext,colback=termbg,colframe=termline,boxrule=0.6pt,arc=3pt,left=8pt,right=8pt,top=5pt,bottom=5pt,
  listing only,before skip=6pt,after skip=8pt,
  overlay={\node[anchor=north east,font=\mono\tiny,text=termline,inner sep=3pt] at (frame.north east) {#1};}}}
\newtcblisting{shell}{termbox/.expanded=bash,listing options={style=term,language=shellx}}
\newtcbinputlisting{\yamlfile}[1]{termbox/.expanded=yaml,listing file=#1,listing options={style=term,language=yamlx,tabsize=2}}

\begin{document}
\thispagestyle{empty}
\begin{tikzpicture}[remember picture,overlay]
  \fill[ink,opacity=0.16] ([xshift=-3.2cm,yshift=-7.2cm]current page.north east) circle (4.6cm);
  \draw[inkline,line width=0.4pt] ([xshift=2.3cm,yshift=-1.6cm]current page.north west) -- ([xshift=-2.3cm,yshift=-1.6cm]current page.north east);
  \node[anchor=north west,inner sep=0] at ([xshift=2.3cm,yshift=-1.0cm]current page.north west) {\eyebrow{Sistemas de Gesti\'on Empresarial \textperiodcentered\ DAM}};
  \node[anchor=north east,inner sep=0] at ([xshift=-2.3cm,yshift=-1.0cm]current page.north east) {\eyebrow{2026/2027}};
  \node[anchor=south west,inner sep=0,text width=15cm] at ([xshift=2.3cm,yshift=7.5cm]current page.south west)
    {{\serif\fontsize{50}{52}\selectfont\spaceskip=0.35em plus 0.1em\color{ink} Odoo 20\\ en una m\'aquina\\ virtual\par}};
  \node[anchor=south west,inner sep=0,text width=15cm] at ([xshift=2.3cm,yshift=5.2cm]current page.south west)
    {{\mono\color{inkdim}Caso pr\'actico 1 \textperiodcentered\ Manual de instalaci\'on y configuraci\'on \textperiodcentered\ 13 oct 2026\par}};
  \draw[inkline,line width=0.4pt] ([xshift=2.3cm,yshift=3.2cm]current page.south west) -- ([xshift=-2.3cm,yshift=3.2cm]current page.south east);
  \node[anchor=north west,inner sep=0] at ([xshift=2.3cm,yshift=3.0cm]current page.south west)
    {\eyebrow{Ubuntu 26.04 \textperiodcentered\ Docker Compose \textperiodcentered\ Odoo 20.0 \textperiodcentered\ PostgreSQL 16}};
\end{tikzpicture}
\newpage
\setcounter{tocdepth}{1}
\renewcommand{\contentsname}{\mono\small\color{inkmuted}\MakeUppercase{\'Indice}}
\tableofcontents
\newpage

\section{Objetivo y entorno}
Instalar Odoo 20 con PostgreSQL 16 mediante Docker Compose dentro de una m\'aquina virtual Ubuntu y acceder a \'el desde el navegador del ordenador anfitri\'on en \texttt{http://localhost:8069}.

\campo{Anfitri\'on}{Windows 11 con VMware Workstation}
\campo{Hipervisor}{VMware Workstation (el enunciado menciona VirtualBox; el reenv\'io de puertos se hizo con el equivalente de VMware)}
\campo{Sistema invitado}{Ubuntu 26.04 LTS (ISO \texttt{ubuntu-26.04-desktop-amd64})}
\campo{Hardware VM}{8 GB de RAM, 4 procesadores, disco de 40 GB}
\campo{Red}{NAT (\texttt{vmnet8}) con reenv\'io de puertos}
\campo{Software}{Docker (paquetes de Ubuntu), Docker Compose v2, imagen \texttt{odoo:20.0}, imagen \texttt{postgres:16}}

\section{Configuraci\'on de red de la VM}
El adaptador de red de la VM est\'a en modo \textbf{NAT}: la VM comparte la IP del anfitri\'on y queda en una red privada, por lo que solo es accesible desde fuera mediante reenv\'io de puertos.

\fig{1.3}{Fig. 1 \textperiodcentered\ Ajustes de la VM: adaptador de red en modo NAT.}

Dentro de la VM, la interfaz \texttt{ens33} recibe por DHCP la IP \texttt{192.168.231.128/24} y la puerta de enlace es \texttt{192.168.231.2}, el router virtual de \texttt{vmnet8}.

\fig{1.1}{Fig. 2 \textperiodcentered\ \cmd{ip -br a} e \cmd{ip route} en la VM.}

En el editor de red virtual de VMware (NAT Settings de \texttt{vmnet8}, subred \texttt{192.168.231.0/24}) se crearon dos reglas de reenv\'io de puertos:

\rowcolors{2}{tint}{paper}
\begin{center}\small
\begin{tabularx}{\textwidth}{@{}l l l X@{}}
\rowcolor{paper}\cab{PUERTO ANFITRI\'ON} & \cab{TIPO} & \cab{DESTINO EN LA VM} & \cab{USO} \\
\arrayrulecolor{inkmuted}\midrule
8069 & TCP & 192.168.231.128:8069 & Odoo \\
2222 & TCP & 192.168.231.128:22 & SSH \\
\arrayrulecolor{inkmuted}\bottomrule
\end{tabularx}
\end{center}
\rowcolors{1}{paper}{paper}

\fig[0.6\textwidth]{1.2}{Fig. 3 \textperiodcentered\ NAT Settings de \texttt{vmnet8} con el reenv\'io de puertos.}

\section{Instalaci\'on de Docker}
Se actualiz\'o el sistema y se instalaron Docker y el plugin de Compose desde los repositorios de Ubuntu (ver el problema 1 en la secci\'on de problemas):

\begin{shell}
$ sudo apt update && sudo apt upgrade -y
$ sudo apt install -y docker.io docker-compose-v2
$ sudo systemctl enable --now docker
$ sudo usermod -aG docker $USER
$ newgrp docker
$ docker --version
$ docker compose version
\end{shell}

Se verific\'o la instalaci\'on con la imagen de prueba:

\begin{shell}
$ docker run --rm hello-world
\end{shell}

\fig{2.2}{Fig. 4 \textperiodcentered\ \cmd{hello-world}: Docker funciona y descarga im\'agenes de Docker Hub.}

\clearpage\section{Fichero docker-compose.yaml}
Se cre\'o el directorio de trabajo \texttt{\~{}/odoo20} y se comprob\'o antes que la etiqueta \texttt{odoo:20.0} existe en el registro:

\fig{2.3}{Fig. 5 \textperiodcentered\ \cmd{docker manifest inspect odoo:20.0}: la imagen existe.}

El fichero se edit\'o con \texttt{nano}. Contenido completo (entregable 1):

\yamlfile{docker-compose.yaml}

Puntos relevantes:
\begin{itemize}
\item El servicio \texttt{web} se conecta a PostgreSQL usando el nombre del servicio (\texttt{HOST: db}); Compose crea una red por defecto con DNS interno.
\item Solo se publica el puerto \texttt{8069} de Odoo; PostgreSQL no se expone fuera de la red de Compose.
\item Los datos persisten en dos vol\'umenes con nombre (\texttt{odoo-db-data} y \texttt{odoo-web-data}).
\item \texttt{--http-interface=0.0.0.0} hace que Odoo escuche en todas las interfaces del contenedor.
\item \texttt{restart: unless-stopped} reinicia los contenedores tras un reinicio de la VM.
\end{itemize}

\fig{2.4}{Fig. 6 \textperiodcentered\ \texttt{docker-compose.yaml} en nano (la l\'inea del volumen de \texttt{db} aparece cortada solo por el ancho de pantalla de nano).}

\section{Arranque y estado de los contenedores}
Se levant\'o la pila con \cmd{docker compose up -d}. Docker descarg\'o las im\'agenes, cre\'o la red \texttt{odoo20\_default}, los dos vol\'umenes y arranc\'o los contenedores \texttt{odoo\_db} y \texttt{odoo\_web}. La salida de \cmd{docker compose ps} muestra ambos en estado \emph{Up} y el puerto \texttt{0.0.0.0:8069->8069/tcp} publicado.

\fig{2.5}{Fig. 7 \textperiodcentered\ Descarga de im\'agenes, creaci\'on de recursos y \cmd{docker compose ps}.}

Tras a\~nadir la l\'inea \texttt{command} al compose (secci\'on de problemas) se volvi\'o a ejecutar \cmd{docker compose up -d} y se comprob\'o desde la propia VM que Odoo responde: \cmd{curl -I http://localhost:8069/web/login} devuelve \texttt{303 SEE OTHER} hacia \texttt{/web/database/selector} y la cabecera \texttt{server: odoo/20.0}. Eso es lo esperado cuando todav\'ia no existe ninguna base de datos.

\fig{3.2}{Fig. 8 \textperiodcentered\ \cmd{docker compose up -d} y \cmd{curl -I} contra Odoo desde la VM.}

\section{Acceso desde el anfitri\'on}
Desde el navegador de Windows (anfitri\'on) se abri\'o \texttt{http://localhost:8069}. VMware reenv\'ia el puerto 8069 a la VM y Odoo muestra el gestor de bases de datos (\texttt{/web/database/selector}), lo que confirma que la cadena anfitri\'on $\to$ NAT $\to$ VM $\to$ contenedor \texttt{web} funciona.

\fig{3.3}{Fig. 9 \textperiodcentered\ Pantalla de inicio de Odoo 20 en el navegador del anfitri\'on, formulario de creaci\'on de base de datos.}

Odoo genera una \emph{master password} cuando el gestor de bases de datos no est\'a protegido; hay que guardarla porque se pide en operaciones posteriores sobre bases de datos (copia, borrado, restauraci\'on).

\section{Creaci\'on de la base de datos e instalaci\'on de CRM}
En el formulario del gestor de bases de datos se rellenaron los siguientes campos y se puls\'o \emph{Create database}:

\campo{Base de datos}{\texttt{MiOdooDatabase}}
\campo{Idioma}{Spanish / Espa\~nol}
\campo{Datos de demostraci\'on}{desactivados}
\campo{Usuario administrador}{el correo y la contrase\~na indicados en el formulario (la contrase\~na aparece oculta)}

\fig{odoo}{Fig. 12 \textperiodcentered\ Formulario de creaci\'on de la base de datos \texttt{MiOdooDatabase}; abajo, los contenedores en ejecuci\'on.}

Al terminar la creaci\'on, Odoo inicia sesi\'on autom\'aticamente con el usuario administrador y abre la aplicaci\'on \emph{Conversaciones} (con el mensaje de bienvenida de OdooBot). En el men\'u de aplicaciones aparece \textbf{CRM} junto a Conversaciones, Calendario, Contactos, Aplicaciones y Ajustes, lo que confirma que la aplicaci\'on est\'a instalada en la base de datos creada y que el objetivo del caso pr\'actico est\'a cumplido.

\fig{odoo2}{Fig. 13 \textperiodcentered\ Odoo 20 con la base de datos creada: el men\'u de aplicaciones incluye CRM.}

\section{Versi\'on de PostgreSQL}
Con la consulta \texttt{SELECT version();} ejecutada dentro del contenedor \texttt{db}:

\begin{shell}
$ docker compose exec db psql -U odoo -d postgres -c "SELECT version();"
\end{shell}

El resultado es \textbf{PostgreSQL 16.15} (Debian 16.15-1.pgdg13+2, x86\_64, 64-bit), coherente con la imagen \texttt{postgres:16}.

\fig{3.1}{Fig. 10 \textperiodcentered\ Versi\'on de PostgreSQL dentro del contenedor \texttt{odoo\_db}.}

\section{Problemas encontrados y soluciones}

\subsection{Docker}
\etiq{Repositorio oficial no disponible}\par\vspace{4pt}
\campo{S\'intoma}{\cmd{sudo apt install docker-ce} y el resto de paquetes oficiales de Docker fallan con \emph{El paquete docker-ce no tiene un candidato para la instalaci\'on}.}
\campo{Causa}{Los paquetes \texttt{docker-ce} no se encuentran para Ubuntu 26.04 (\emph{resolute}) en los repositorios configurados. La causa exacta (repositorio de Docker sin a\~nadir o sin soporte para esa versi\'on) \nv.}
\campo{Soluci\'on}{Usar los paquetes de los repositorios de Ubuntu: \cmd{sudo apt install -y docker.io docker-compose-v2}. Antes se intent\'o borrar \texttt{/etc/apt/sources.list.d/docker.list}, que no exist\'ia (\emph{No existe el archivo}), sin efecto.}

\fig{2.1}{Fig. 11 \textperiodcentered\ Error de instalaci\'on de \texttt{docker-ce} y cambio a \texttt{docker.io}.}

\subsection{Permisos}
\etiq{Docker sin sudo}\par\vspace{4pt}
\campo{Soluci\'on}{A\~nadir el usuario al grupo \texttt{docker} (\cmd{sudo usermod -aG docker \$USER}) y activar el grupo en la sesi\'on actual con \cmd{newgrp docker}, sin cerrar sesi\'on.}

\subsection{Odoo}
\etiq{Escucha en todas las interfaces}\par\vspace{4pt}
\campo{Cambio}{La primera versi\'on del compose (31 l\'ineas) no inclu\'ia \texttt{command}. Despu\'es se a\~nadi\'o \texttt{command: ["odoo", "--http-interface=0.0.0.0"]} (32 l\'ineas) y se recre\'o la pila.}
\campo{Motivo}{Garantizar que Odoo acepta conexiones desde fuera del contenedor, necesarias para el puerto publicado y el reenv\'io de VMware. Si fue estrictamente necesario con la imagen \texttt{odoo:20.0} \nv.}

\end{document}
